% Skeleton draft. Uses plain article + manual formatting as a stand-in for the
% official venue class, since the exact NeurIPS 2026 workshop .sty/.cls was not
% fetched in this pass (OpenReview submissions opened 2026-07-15, the day this
% skeleton was written -- the template link may not even be live yet). BEFORE
% SUBMISSION: download the official template from
% https://efficient-ondevice-ai-agents.github.io/ (or the OpenReview venue
% page), drop the .sty/.cls file into this directory, and swap the
% documentclass/preamble below for theirs. Everything else (content, section
% structure, citations) should carry over with minimal changes.
%
% Target: short paper, 4 pages + unlimited references (docs/workshop_info.md).
\documentclass[10pt,twocolumn]{article}

\usepackage[margin=0.75in,columnsep=0.25in]{geometry}
\usepackage{times}
\usepackage[T1]{fontenc}
\usepackage{natbib}
\usepackage{graphicx}
\usepackage{booktabs}
\usepackage{amsmath}
\usepackage{hyperref}
\hypersetup{colorlinks=true, linkcolor=blue, citecolor=blue, urlcolor=blue}

\title{\vspace{-1em}Extending AlphaMaze's Recipe to Gemma 4:\\
First Small-Model Results on GridRoute}
\author{Vedang Pandey \\ \small{Manuscript in preparation -- Efficient and On-Device AI Agents Workshop @ NeurIPS 2026}}
\date{}

\begin{document}
\maketitle

\begin{abstract}
GridRoute~\citep{gridroute} has, to our knowledge, only ever been evaluated with models of 7B
parameters or larger (GPT-4 Turbo, Qwen2.5-7B/72B) -- no result for a small, on-device-feasible
model (under 8B parameters) appears in the benchmark's own paper or in anything citing it.
Separately, AlphaMaze~\citep{alphamaze} showed that SFT followed by GRPO can take a
1.5B-parameter model from near-zero to 93\% accuracy on its own token-formatted maze benchmark,
MazeBench, but never tested a different model family or a natural-language task format. We
extend AlphaMaze's training recipe to Gemma 4, targeting strong performance on GridRoute (natural
-language grid navigation, $5\times5$) without regressing MazeBench (token format), and report the
first small-model numbers on GridRoute we are aware of. We try three training recipes -- GridRoute
-only, a naive mix of both formats, and a mix with an added cross-format-consistency reward
adapted from cross-lingual consistency-reward RL~\citep{elhady2026} -- and report whichever
performs best on both benchmarks jointly, not a specific technique as the point of the paper.
\textbf{[Placeholder: results pending training runs, see \S\ref{sec:results}.]} A harness built to
avoid measurement artifacts we demonstrate are real in our own development history (a single
model's apparent MazeBench accuracy swung from 0\% to 99\% across successive versions of our own
early evaluation code), all code, and a Kaggle-runnable training notebook are released alongside
this draft.
\end{abstract}

\section{Introduction}

Small, on-device language models are increasingly fine-tuned for narrow agentic skills rather than
general capability, and spatial navigation is a natural test case: cheap to generate supervision
for, easy to verify programmatically, and directly relevant to on-device agents that plan routes
or follow spatial instructions under tight compute budgets. Two independent lines of prior work
set up this paper's question but do not answer it. First, GridRoute~\citep{gridroute} -- a
benchmark for natural-language grid navigation with cardinal movement -- has only ever been
evaluated, in its own paper and in every follow-up we found, with models of 7B parameters or
larger. Whether a genuinely small model (under 8B, plausibly on-device-deployable) can be made to
perform well at this task at all is, to our knowledge, simply unmeasured. Second,
AlphaMaze~\citep{alphamaze} demonstrated that supervised fine-tuning followed by Group Relative
Policy Optimization (GRPO) can take a 1.5B model from essentially no maze-solving ability to 93\%
accuracy -- but only for one model family, in one token-based maze representation the model was
specifically trained on, and never tested against a natural-language task format or a different
base model.

This paper closes the gap between these two facts directly: we take Gemma 4 -- a current-generation
small on-device model family -- and apply AlphaMaze's SFT+GRPO recipe with the specific goal of
making it as good as possible at GridRoute, while checking that it does not lose ground on
MazeBench in the process. We do not propose a new training algorithm. We try three plausible
recipes (single-format GRPO, naive mixed-format GRPO, and mixed-format GRPO with a cross-format
-consistency reward adapted from recent cross-lingual work~\citep{elhady2026}) and report whichever
works, including a negative result if none meaningfully improves on the simplest baseline. The
contribution is empirical and methodological, not algorithmic: a first measurement where none
existed, produced by a harness we show is actually necessary.

\textbf{Why the harness matters, concretely.} In the course of this project's own development, a
single fixed model's measured MazeBench accuracy took values of 0\%, 70\%, 88\%, and 99\% across
successive versions of an evaluation harness that varied decoding temperature, thinking-token
budget, and answer-extraction logic without controlling for them independently (\S\ref{sec:setup}).
Separately, an early version of our own harness scored MazeBench by exact string match against one
stored reference solution, when the benchmark's real evaluation methodology (which we adopt
directly, see \S\ref{sec:setup}) simulates a candidate's moves against the maze's actual wall
structure and accepts any sequence that reaches the target. Both are the kind of error that,
uncaught, silently produces a wrong headline number. We report our numbers as an existence proof
that they were checked for exactly this.

\textbf{Contributions.} (1) The first small-model (under 8B) results on GridRoute we are aware of.
(2) An extension of AlphaMaze's SFT+GRPO recipe to Gemma 4 and to a natural-language task format,
evaluated jointly with the original token-format benchmark, with a best-effort search over three
plausible training recipes. (3) An evaluation harness -- clean-final-answer-only scoring, verified
generous thinking budgets, and direct reuse of AlphaMaze's own scoring code rather than a
reconstruction -- built specifically to avoid measurement artifacts we demonstrate are real, plus
a Kaggle-runnable, open-source implementation of the entire pipeline.

\section{Related Work}

\textbf{GridRoute and small models.} GridRoute~\citep{gridroute} defines rectangular-obstacle grid
navigation with cardinal movement and reports GPT-4 Turbo and Qwen2.5-7B/72B under vanilla,
chain-of-thought, and algorithm-of-thought prompting. No model below 7B parameters appears in that
paper or, to our knowledge, in any later work using the benchmark. This is the specific,
literature-verified gap this paper addresses -- not a novelty argument about method, but an
absence-of-evidence argument about scale.

\textbf{Fine-tuning small models for spatial reasoning.} AlphaMaze~\citep{alphamaze} is the direct
methodological precedent: SFT+GRPO on DeepSeek-R1-Distill-Qwen-1.5B, taking token-format maze
accuracy on MazeBench from near zero to 93\%. It does not test a different model family, a
different task format, or GridRoute specifically. \citet{jietal2025} show GRPO generalizes better
than SFT for spatial-reasoning out-of-distribution generalization on a vision-language model, with
OOD meaning reworded queries within one format and modality -- a different setting from testing a
new model on a new format entirely, which is what this paper does.

\textbf{Cross-format training recipes.} Among the three recipes we try, one (\S\ref{sec:method})
augments GRPO with a reward term that scores agreement between a model's own completions on two
representations of the same problem. This mechanism is not new: \citet{elhady2026} introduce it for
cross-\emph{lingual} math reasoning (MGSM) and report real gains there. We apply the same
mechanism to cross-\emph{format} (not cross-lingual) spatial reasoning as one candidate recipe
among the three we test -- this is a transfer-to-a-new-domain application, reported as one line in
a results table, not this paper's contribution. Two superficially similar lines of work,
GRPO-CARE~\citep{grpocare2026} and Faithful GRPO~\citep{faithfulgrpo2026}, define "consistency" as
within-response coherence between a model's reasoning trace and its own final answer -- a
different, unrelated notion. \citet{beyondspecialization2026} show that diverse training data
improves cross-environment transfer for classical (non-LLM) reinforcement learning navigation
policies; our naive mixed-format recipe is the direct LLM/GRPO analogue of that finding.

\section{Method}
\label{sec:method}

\subsection{Two Task Formats}

MazeBench's token format (real published data, \texttt{Menlo/Maze-Bench-v0.2}) is used as-is, no
conversion. For GridRoute, we use natural-language coordinate instructions (start, goal, and
obstacle coordinates, asking for a Python list of \texttt{(row, col)} tuples as the answer) at
grid size $5\times5$, matched to MazeBench's own grid size for a fair joint evaluation and chosen
as an appropriately modest difficulty level for small models -- larger grids ($10\times10$ and up)
are left as future work pending how well $5\times5$ goes. For the mixed-format and consistency
recipes below, we additionally render GridRoute grids in our own token-based encoding (reusing
AlphaMaze's published token vocabulary with our own fully-specified per-cell layout, rather than a
reverse-engineered clone of their undocumented exact grammar -- see the released code for the
exact specification), which lets the same underlying grid be presented in both formats during
training.

\subsection{Training Recipes}

All recipes apply an SFT warm-start (teaching the exact answer-reporting convention scored at
evaluation time) followed by GRPO, on Gemma 4 E2B (and E4B, VRAM permitting).

\textbf{Single-format}: GRPO on GridRoute NL only.
\textbf{Mixed-format}: GRPO on GridRoute NL and our token-maze encoding, interleaved over the same
underlying grids, each scored by its own per-format reward.
\textbf{Consistency-reward}: identical mixed-format data, but the reward additionally greedy-
decodes the model's current completion on the paired prompt in the other format and adds a bonus
when both agree on validity or optimality for the same underlying problem -- requiring
independent correctness on both, not mere matching output, so two confidently-wrong-but-matching
completions earn nothing extra.

We report whichever recipe performs best jointly across both benchmarks (\S\ref{sec:results}),
not a comparison designed to favor any one of them.

\subsection{Evaluation Methodology}
\label{sec:harness}

Two choices here matter enough to state explicitly, motivated directly by errors caught during
this project's own development (\S\ref{sec:results-caveat}):

\textbf{Score only the model's own reported final answer.} A model's answer is whatever it states
after its thinking closes, or after an explicit "final answer" instruction we give it for the NL
format; if neither ever appears, or generation is truncated before either does, that is scored as
a failure. Coordinates or moves are never extracted from a thinking trace the model did not itself
present as its answer -- a chain-of-thought model routinely mentions and discards candidate moves
before its real answer, and naive extraction over the full output text conflates the two.

\textbf{Generous, and verified, thinking budgets.} Truncating a reasoning model's thinking mid
-thought measurably changes its apparent accuracy (\S\ref{sec:results-caveat}); we use thinking
budgets large enough that generations reliably complete, confirmed empirically rather than assumed.

\section{Experimental Setup}
\label{sec:setup}

\textbf{Models.} Gemma 4 E2B (primary) and E4B (if its LoRA memory footprint fits the training
hardware, confirmed empirically rather than assumed), via bitsandbytes 4-bit + \texttt{peft} LoRA
(\texttt{peft}$\,\geq\,$0.19.0; Gemma 4's \texttt{Gemma4ClippableLinear} layers need this version to
be a recognized target module type at all). AlphaMaze-v0.2-1.5B as the MazeBench-side reference
point. DeepSeek-R1-Distill-Qwen-1.5B / SmolLM2-1.7B / Qwen2.5-3B as cheap secondary comparisons, via
the same bitsandbytes 4-bit + \texttt{peft} LoRA path.

\textbf{Benchmarks.} MazeBench-v0.2 (100 held-out mazes), scored with AlphaMaze's own released
evaluation code (candidate moves simulated against the maze's actual wall structure, parsed
directly out of the prompt text; any sequence reaching the target counts as correct) rather than a
reimplementation. GridRoute-style NL navigation, $5\times5$, seeded generation; valid-path rate
and optimal-path rate against a Dijkstra reference, as in \citet{gridroute}.

\textbf{Reproducibility.} Fixed seed (42) throughout. All evaluation decoding is greedy
(temperature 0). Code, including the full training/evaluation harness and a Kaggle notebook
reproducing every run in this paper end-to-end, is released at the project repository.

\subsection{A Caveat From Our Own Development History}
\label{sec:results-caveat}

Before arriving at the methodology in \S\ref{sec:harness}, an earlier version of our own harness
measured the identical AlphaMaze checkpoint at 99\%, 70\%, 88\%, and 0\% on MazeBench across four
separate runs within roughly 24 hours, driven by uncontrolled changes to decoding temperature,
thinking-token budget, and answer-extraction logic, plus (separately) scoring by exact string match
against one stored reference solution rather than simulating candidate moves against the maze's
actual structure -- the latter approach can fail a valid-but-different solution that the
benchmark's own methodology would accept. We report this not as an aside but as direct evidence
for why \S\ref{sec:harness}'s choices are load-bearing: a small-model spatial-reasoning number
produced without this care is not obviously trustworthy, on this evidence.

\section{Results}
\label{sec:results}

\textbf{Results are pending as of this draft.} Baselines (Gemma 4 E2B/E4B and AlphaMaze,
untrained, on both benchmarks) and the three training recipes had not yet been executed on GPU
hardware at the time this skeleton was written. Table~\ref{tab:main} shows the structure the final
results table will follow.

\begin{table}[h]
\centering
\small
\begin{tabular}{lcc}
\toprule
Condition & MazeBench (token) & GridRoute (NL) \\
\midrule
AlphaMaze ckpt. (untrained) & TBD & TBD \\
Gemma 4 E2B (untrained) & TBD & TBD \\
Gemma 4 E2B, single-format & TBD & TBD \\
Gemma 4 E2B, mixed-format & TBD & TBD \\
Gemma 4 E2B, consistency & TBD & TBD \\
\bottomrule
\end{tabular}
\caption{Planned central results table (accuracy / valid-path rate by condition and benchmark).
To be filled in after training runs; Gemma 4 E4B rows added if it fits the training hardware.}
\label{tab:main}
\end{table}

\section{Discussion and Limitations}

\textbf{This is a measurement and engineering contribution, not a new algorithm.} The consistency
-reward recipe is Elhady et al.'s~\citep{elhady2026} mechanism, applied to a new domain as one of
three recipes tried, not proposed here as a method. If it happens to win, that is reported as "an
existing technique transfers here too," not as a contribution in its own right.

\textbf{Our token-maze format is not AlphaMaze's exact format} (\S\ref{sec:method}) -- if the
mixed or consistency recipes show no effect, this is a confound to check before concluding the
recipe itself doesn't help.

\textbf{A lightweight failure-category breakdown} (no output, start/end mismatch, obstacle
collision, out-of-bounds, invalid step) accompanies the main table; a fuller per-recipe failure
analysis is left as future work given the time available for this submission.

\section{Conclusion}

We report, to our knowledge, the first results for a small (under 8B parameter) on-device language
model on GridRoute, by extending AlphaMaze's SFT+GRPO recipe to Gemma 4 and testing three training
recipes for joint performance on GridRoute and MazeBench. We make no algorithmic claim; the
contribution is an absent measurement, taken carefully, with the harness choices motivated
directly by errors we demonstrate are real in our own development history. Final results,
including negative ones, will determine which (if any) of the three recipes is worth
recommending.

\bibliographystyle{plainnat}
\bibliography{references}

\end{document}
